%% Manual for inspirehep.sty.  Build with:  make doc
\documentclass[a4paper,11pt]{article}
\usepackage[T1]{fontenc}
\usepackage[margin=1in]{geometry}
\usepackage[hidelinks]{hyperref}
\usepackage{booktabs}
\usepackage{fancyvrb}
\usepackage{inspirehep}

% \BibTeX is not in the kernel; define it rather than depend on a package.
\providecommand{\BibTeX}{B\textsc{ib}\TeX}
\newcommand{\pkg}[1]{\textsf{#1}}
\newcommand{\opt}[1]{\texttt{#1}}
\setlength{\parindent}{0pt}
\setlength{\parskip}{0.6em}

\title{The \pkg{inspirehep} package\\[0.3em]
  \large Live INSPIRE-HEP data in a \LaTeX{} document}
\author{Lawrence Lee}
\date{Version 2.0 \quad 2026/08/29}

\begin{document}
\maketitle

\begin{abstract}
\noindent
\pkg{inspirehep} puts data from \href{https://inspirehep.net}{INSPIRE-HEP} into
a \LaTeX{} document: hand a command a record id and it typesets the title, the
citation count, the full reference or the \BibTeX{} entry; hand it a person's id
and it gives their publication count, citation count or $h$-index, or a plot of
either against time. The numbers live in a generated file, so the document
always compiles --- offline, on a stranger's machine, on Overleaf --- whether or
not it can reach the network.
\end{abstract}

\section{Installation}
Put \texttt{inspirehep.sty} where \LaTeX{} will find it: beside your document,
or anywhere on \texttt{TEXINPUTS}. It requires only \pkg{xparse} and
\pkg{l3keys2e}, both part of a current \LaTeX{}; \pkg{pgfplots} is needed only
if you plot, and is loaded only if you ask for it with the \opt{plots} option.

\section{Getting the data}
The figures are read from \texttt{inspirehep-data.tex}. Two things can write
it, and they write the same file.

\textbf{The package itself.} With unrestricted shell escape it refreshes at the
end of a run and asks for a rerun, exactly like a cross-reference:
\begin{Verbatim}[frame=single,fontsize=\small]
pdflatex -shell-escape cv.tex
pdflatex -shell-escape cv.tex
\end{Verbatim}
After that it only refetches when there is a reason: a record you have just
added, or data older than \opt{maxage}.

\textbf{Or the helper.} \texttt{inspirehep-fetch.py} does the same job from
outside the compile, needs only the Python standard library, and reads your
sources to work out what to ask for:
\begin{Verbatim}[frame=single,fontsize=\small]
python3 inspirehep-fetch.py cv.tex
\end{Verbatim}

\subsection{Overleaf}
Overleaf disables shell escape and its compile containers have no network, so
nothing can fetch there --- but a document still typesets correct figures,
because they are only an \verb|\input| file. Upload \texttt{inspirehep.sty},
then add \texttt{inspirehep-data.tex} through \emph{Add file $\to$ From external
URL}, pointing at the raw copy in your repository. Overleaf gives linked files
a \emph{Refresh} button, so new figures are one click away. Note that Overleaf's
Git sync does not fetch submodules: include the \texttt{.sty} as a file.

\section{Records}
\begin{Verbatim}[frame=single,fontsize=\small]
\inspirepub{2642414}                  \inspirekey{2642414}
\inspirepub[year]{2642414}            \inspirecite{2642414}
\inspirepub[ref]{2642414}             \inspiretitle{2642414}
\inspirepub[cites=false]{2642414}     \inspirecites{2642414}
\inspirepub[title={My words}]{2642414}  \inspireref{2642414}
\end{Verbatim}
The id is the number in the INSPIRE URL,
\texttt{inspirehep.net/literature/\textbf{2642414}}.

References and \BibTeX{} come from INSPIRE's own renderers rather than being
assembled from the metadata, so they match what everyone else in the field
cites, errata included.

\section{People}
Every author command takes the id it is asking about, so one document can
discuss several people. The id is the number in a profile URL,
\texttt{inspirehep.net/authors/\textbf{1071846}}, or a BAI such as
\texttt{J.Smith.1}.
\begin{Verbatim}[frame=single,fontsize=\small]
\inspirepapers{1071846}      \inspirecitations[round=1000]{1071846}
\inspirehindex{1071846}      \inspireauthorstat{1071846}{citations}
\end{Verbatim}
\opt{round} rounds \emph{down}, so ``over $N$'' stays true as the real figure
grows.

\section{Plots}
With the \opt{plots} option, \verb|\inspireplot{<id>}| draws citations per year
for a paper and \verb|\inspireauthorplot{<id>}{papers|\texttt{|}\verb|citations}|
the same for a person. The default is one line with axis lines only where they
carry information; \verb|\inspireplotstyle| takes any \pkg{pgfplots} axis keys.

\section{Options}
Every option below is both a package option and a per-call option, and
\verb|\inspiresetup{...}| changes any of them mid-document.

\begin{center}\small
\begin{tabular}{@{}lll@{}}
\toprule
Option & Default & Meaning \\
\midrule
\opt{cites}    & \opt{true}  & show the citation count \\
\opt{ref}      & \opt{false} & the full reference instead of the title \\
\opt{year}     & \opt{false} & append the publication year \\
\opt{link}     & \opt{true}  & hyperlink the title \\
\opt{title}    & ---         & override the fetched title \\
\opt{round}    & \opt{1}     & round a figure down to a multiple of this \\
\opt{sep}      & \opt{comma} & \opt{comma}, \opt{period}, \opt{space}, \opt{thin}, \opt{underscore}, \opt{none} \\
\opt{sepstring}& ---         & any other separator \\
\opt{style}    & \opt{latex-eu} & reference format: \opt{latex-eu} or \opt{latex-us} \\
\opt{plots}    & \opt{false} & load \pkg{pgfplots} \\
\opt{data}     & \opt{inspirehep-data} & basename of the generated file \\
\opt{bib}      & \opt{inspirehep-refs} & basename of the generated \texttt{.bib} \\
\opt{maxage}   & \opt{120}   & days before the data is called stale \\
\opt{mincites} & \opt{1}     & counts below this print nothing \\
\opt{fetch}    & \opt{auto}  & \opt{auto}, \opt{on} or \opt{off} \\
\bottomrule
\end{tabular}
\end{center}

\section{Changing how things look}
\begin{Verbatim}[frame=single,fontsize=\small]
\renewcommand{\inspiretitleformat}[1]{\textbf{#1}}
\renewcommand{\inspirecitestext}[1]{#1~cites}
\renewcommand{\inspirecitesformat}[1]{\nobreakspace{\small[#1]}}
\renewcommand{\inspireyearformat}[1]{\hfill #1}
\renewcommand{\inspireplotstyle}{ymajorgrids}
\end{Verbatim}

\section{Limits}
Fetching from the package needs \texttt{curl} and unrestricted shell escape;
the helper is the way round both. Counts are whatever INSPIRE reports,
including self-citations.

\section{Licence}
LPPL 1.3c or later. Maintained by Lawrence Lee;
\url{https://github.com/lawrenceleejr/inspirehep-latex}.

\end{document}
